% GENERATED FILE. Edit canonical lessons, then run npm run generate:books.
% canonical-source-hash: fnv1a64:52113d9d418d876b
% canonical-lessons: RU-C301-katastrofa, RU-C301-vzryv, RU-C301-vzryvatsya, RU-C301-obrushitsya, RU-C301-panika

\chapter{Disaster}
\label{ch:ru-disaster}
\hlchaptermodality{301}

\begin{chapteropening}
\textbf{By the end of this chapter:} \emph{I can talk about a disaster, an explosion, something blowing up, a building collapsing and panic.}

Complete the last lesson of chapter 301: I can talk about a disaster, an explosion, something blowing up, a building collapsing and panic.
\end{chapteropening}

\section[\texorpdfstring{\ru{катастрофа} (katastrófa)}{катастрофа (katastrófa)}]{\ru{катастрофа} (katastrófa) — a disaster, a catastrophe; a crash}

\label{lesson:RU-C301-katastrofa}



\begin{quote}
Before the new one: say the Russian for a drought, then the Russian for to flood.
\end{quote}

\subsection*{You'll want to know: \ru{катастрофа}}
\textbf{\ru{катастрофа}} (\emph{katastrófa}) — \textquotedblleft{}a disaster, a catastrophe; a crash\textquotedblright{}.

\begin{quote}
\ru{Это} \ru{была} \ru{настоящая} \ru{катастрофа}.
\end{quote}

\begin{quote}
\emph{Éto bylá nastoyáshchaya katastrófa.}
\end{quote}

\begin{quote}
It was a real disaster.
\end{quote}

Also a crash: \emph{\ru{авиакатастрофа}} — a plane crash.

\begin{cousinweb}
\emph{\ru{катастрофа}} goes back to Greek \emph{katastroph\'{\={e}}} \textquotedblleft{}an overturning\textquotedblright{}, from \emph{katastréphō} \textquotedblleft{}I overturn\textquotedblright{}, built on \emph{katá} \textquotedblleft{}down\textquotedblright{} and \emph{stréphō} \textquotedblleft{}I turn\textquotedblright{} (EtymDB).
\end{cousinweb}

\begin{grammarlens}[title={What you've built}]
One more word for telling, explaining and describing at B1.
\end{grammarlens}

\subsection*{Your turn}
\begin{itemize}
\raggedright
  \item \emph{Say it:} \emph{katastrófa}
  \item \emph{Say it:} \emph{katastrófa}, once more
  \item \emph{Say it:} a sentence of your own with \emph{katastrófa}
  \item \emph{Recall:} say the Russian for a drought, then the Russian for to flood, then say \emph{katastrófa} again
\end{itemize}

\begin{tcolorbox}[breakable,colback=teal!4,colframe=teal!35!black,title={Before you move on}]
What does \ru{катастрофа} mean? (\textquotedblleft{}A disaster, a catastrophe; a crash\textquotedblright{}.) Say it once more.
\end{tcolorbox}
\section[\texorpdfstring{\ru{взрыв} (vzryv)}{взрыв (vzryv)}]{\ru{взрыв} (vzryv) — an explosion}

\label{lesson:RU-C301-vzryv}



\begin{quote}
Before the new one: say the Russian for to flood, then the Russian for a disaster.
\end{quote}

\subsection*{You'll want to know: \ru{взрыв}}
\textbf{\ru{взрыв}} (\emph{vzryv}) — \textquotedblleft{}an explosion\textquotedblright{}.

\begin{quote}
\ru{Мы} \ru{услышали} \ru{громкий} \ru{взрыв}.
\end{quote}

\begin{quote}
\emph{My uslýshali grómkiy vzryv.}
\end{quote}

\begin{quote}
We heard a loud explosion.
\end{quote}

Masculine.

\begin{cousinweb}
\emph{\ru{взрыв}} is built on \emph{\ru{взрывать}} \textquotedblleft{}to blow up\textquotedblright{}, which is usually traced to \emph{\ru{вз}-} \textquotedblleft{}up\textquotedblright{} and \emph{\ru{рыть}} \textquotedblleft{}to dig\textquotedblright{}: to throw the earth up. EtymDB lists no parent.
\end{cousinweb}

\begin{grammarlens}[title={What you've built}]
One more word for telling, explaining and describing at B1.
\end{grammarlens}

\subsection*{Your turn}
\begin{itemize}
\raggedright
  \item \emph{Say it:} \emph{vzryv}
  \item \emph{Say it:} \emph{vzryv}, once more
  \item \emph{Say it:} a sentence of your own with \emph{vzryv}
  \item \emph{Recall:} say the Russian for to flood, then the Russian for a disaster, then say \emph{vzryv} again
\end{itemize}

\begin{tcolorbox}[breakable,colback=teal!4,colframe=teal!35!black,title={Before you move on}]
What does \ru{взрыв} mean? (\textquotedblleft{}An explosion\textquotedblright{}.) Say it once more.
\end{tcolorbox}
\section[\texorpdfstring{\ru{взрываться} (vzryvátsya)}{взрываться (vzryvátsya)}]{\ru{взрываться} (vzryvátsya) — to explode, to blow up}

\label{lesson:RU-C301-vzryvatsya}



\begin{quote}
Before the new one: say the Russian for a disaster, then the Russian for an explosion.
\end{quote}

\subsection*{You'll want to know: \ru{взрываться}}
\textbf{\ru{взрываться}} (\emph{vzryvátsya}) — \textquotedblleft{}to explode, to blow up\textquotedblright{}.

\begin{quote}
\ru{На} \ru{кухне} \ru{взорвался} \ru{газ}.
\end{quote}

\begin{quote}
\emph{Na kúkhne vzorválsya gaz.}
\end{quote}

\begin{quote}
The gas exploded in the kitchen.
\end{quote}

Imperfective; perfective \emph{\ru{взорваться}} (\emph{vzorvátsya}), \emph{\ru{взорвётся}} (\emph{vzorvyótsya}), on another stem.

\begin{cousinweb}
\emph{\ru{взрываться}} is the reflexive of \emph{\ru{взрывать}} \textquotedblleft{}to blow up\textquotedblright{}, which is usually traced to \emph{\ru{вз}-} \textquotedblleft{}up\textquotedblright{} and \emph{\ru{рыть}} \textquotedblleft{}to dig\textquotedblright{}; EtymDB lists no parent.
\end{cousinweb}

\begin{grammarlens}[title={What you've built}]
One more word for telling, explaining and describing at B1.
\end{grammarlens}

\subsection*{Your turn}
\begin{itemize}
\raggedright
  \item \emph{Say it:} \emph{vzryvátsya}
  \item \emph{Say it:} \emph{vzryvátsya}, once more
  \item \emph{Say it:} a sentence of your own with \emph{vzryvátsya}
  \item \emph{Recall:} say the Russian for a disaster, then the Russian for an explosion, then say \emph{vzryvátsya} again
\end{itemize}

\begin{tcolorbox}[breakable,colback=teal!4,colframe=teal!35!black,title={Before you move on}]
What does \ru{взрываться} mean? (\textquotedblleft{}To explode, to blow up\textquotedblright{}.) Say it once more.
\end{tcolorbox}
\section[\texorpdfstring{\ru{обрушиться} (obrúshitsya)}{обрушиться (obrúshitsya)}]{\ru{обрушиться} (obrúshitsya) — to collapse, to cave in}

\label{lesson:RU-C301-obrushitsya}



\begin{quote}
Before the new one: say the Russian for an explosion, then the Russian for to explode.
\end{quote}

\subsection*{You'll want to know: \ru{обрушиться}}
\textbf{\ru{обрушиться}} (\emph{obrúshitsya}) — \textquotedblleft{}to collapse, to cave in\textquotedblright{}.

\begin{quote}
\ru{Старый} \ru{мост} \ru{обрушился}.
\end{quote}

\begin{quote}
\emph{Stáryy most obrúshilsya.}
\end{quote}

\begin{quote}
The old bridge collapsed.
\end{quote}

Perfective; imperfective \emph{\ru{обрушиваться}} (\emph{obrúshivatsya}).

\begin{cousinweb}
\emph{\ru{обрушиться}} is built on \emph{\ru{об}-} and \emph{\ru{рушить}} \textquotedblleft{}to pull down\textquotedblright{}, an old East Slavic verb; EtymDB lists \emph{\ru{рушить}} with no parent.
\end{cousinweb}

\begin{grammarlens}[title={What you've built}]
One more word for telling, explaining and describing at B1.
\end{grammarlens}

\subsection*{Your turn}
\begin{itemize}
\raggedright
  \item \emph{Say it:} \emph{obrúshitsya}
  \item \emph{Say it:} \emph{obrúshitsya}, once more
  \item \emph{Say it:} a sentence of your own with \emph{obrúshitsya}
  \item \emph{Recall:} say the Russian for an explosion, then the Russian for to explode, then say \emph{obrúshitsya} again
\end{itemize}

\begin{tcolorbox}[breakable,colback=teal!4,colframe=teal!35!black,title={Before you move on}]
What does \ru{обрушиться} mean? (\textquotedblleft{}To collapse, to cave in\textquotedblright{}.) Say it once more.
\end{tcolorbox}
\section[\texorpdfstring{\ru{паника} (pánika)}{паника (pánika)}]{\ru{паника} (pánika) — panic}

\label{lesson:RU-C301-panika}



\begin{quote}
Before the new one: say the Russian for to explode, then the Russian for to collapse.
\end{quote}

\subsection*{You'll want to know: \ru{паника}}
\textbf{\ru{паника}} (\emph{pánika}) — \textquotedblleft{}panic\textquotedblright{}.

\begin{quote}
\ru{Люди} \ru{бежали} \ru{в} \ru{панике}.
\end{quote}

\begin{quote}
\emph{Lyúdi bezháli v pánike.}
\end{quote}

\begin{quote}
People were running in panic.
\end{quote}

Feminine. \emph{\ru{Не} \ru{паникуй}!} — don't panic!

\begin{cousinweb}
\emph{\ru{паника}} is usually traced through French \emph{panique} to Greek \emph{panikós} \textquotedblleft{}of Pan\textquotedblright{}, the god said to bring sudden fear; EtymDB lists no parent.
\end{cousinweb}

\begin{grammarlens}[title={What you've built}]
One more word for telling, explaining and describing at B1.
\end{grammarlens}

\subsection*{Your turn}
\begin{itemize}
\raggedright
  \item \emph{Say it:} \emph{pánika}
  \item \emph{Say it:} \emph{pánika}, once more
  \item \emph{Say it:} a sentence of your own with \emph{pánika}
  \item \emph{Recall:} say the Russian for to explode, then the Russian for to collapse, then say \emph{pánika} again
\end{itemize}

\begin{tcolorbox}[breakable,colback=teal!4,colframe=teal!35!black,title={Before you move on}]
What does \ru{паника} mean? (\textquotedblleft{}Panic\textquotedblright{}.) Say it once more.
\end{tcolorbox}
